\subsection{Diagonalization}
\paragraph*{Definition}
A matrix $A \in \mathbb{R}^{n \times n}$ is diagonalizable if there exists an invertible matrix $P$ and a diagonal matrix $D$ such that
\begin{equation*}
    A = PDP^{-1} 
\end{equation*}

Diagonal and triangular matricies have their eigenvalues on the main diagonal.

\paragraph*{Theorem}
An $n \times n$ matrix $A$ is said to be diagonalizable if and only if $P^{-1}AP = D$ whe $P$ is some invertible matrix and $D$ is a diagonal matrix. 
In this case, the diagonal entries of $D$ are the eigenvalues of $A$ and the columns of $P$ are the corresponding eigenvectors of $A$.

\paragraph*{Properties}
\begin{itemize}
    \item $C = $ diag$(\lambda_1, \lambda_2, \cdots, \lambda_n)$ is diagonal $\implies$ $C$ is diagonalizable
    \item $CD = $ diag$(c_1d_1, c_2d_2, \cdots, c_nd_n) = CD$ is diagonal $\implies$ $CD$ is diagonalizable
\end{itemize}

\begin{equation*}
    A = \begin{bmatrix}
        \lambda_1 & 0 & 0 & 0\\
        0 & \lambda_2 & 0 & 0 \\
        0 & 0 & \ddots & \vdots \\
        0 & 0 & \cdots & \lambda_n
    \end{bmatrix}
\end{equation*}

\paragraph*{Theorem}
An $n \times n$ matrix $A$ has n eigenvalues
\begin{itemize}
    \item Not necessarily distinct
    \item Not necessarily real
    \item $\det(A) = \lambda_1 \lambda_2 \cdots \lambda_n$
    \item $\lambda = 0 \implies \det(A) = 0 \implies A$ is not invertible
\end{itemize}

\subsubsection*{Similarity}
\paragraph*{Definition}
Two $n \times n$ matrices $A$ and $B$ are similar $(A \sim B)$ if there exists an invertible matrix $P$ such that $B = P^{-1}AP$. 
Similar matrices have the same characteristic polynomial and therefore the same eigenvalues.

\paragraph*{How do we find $P$ and $D$?}
We want to find $P$ such that $P^{-1}AP = D$ where $D$ is diagonal.
\begin{align*}
    P^{-1}&AP = D \implies AP = PD\\
    P = &[\vec{v_1} \vec{v_2} \cdots \vec{v_n}] \implies A[\vec{v_1} \vec{v_2} \cdots \vec{v_n}] = [\vec{v_1} \vec{v_2} \cdots \vec{v_n}]\begin{bmatrix}
        \lambda_1 & 0 & 0 & 0\\
        0 & \lambda_2 & 0 & 0 \\
        0 & 0 & \ddots & \vdots \\
        0 & 0 & \cdots & \lambda_n
    \end{bmatrix} \\
    \implies &[A\vec{v_1} A\vec{v_2} \cdots A\vec{v_n}] = [\lambda_1 \vec{v_1} \lambda_2 \vec{v_2} \cdots \lambda_n \vec{v_n}]\\
    \implies &A\vec{v_i} = \lambda_i \vec{v_i} \text{ for } i = 1, 2, \cdots, n\\
    \therefore \quad &\vec{v_1}, \vec{v_2}, \cdots, \vec{v_n} \text{ are eigenvectors of } A \text{ corresponding to } \lambda_1, \lambda_2, \cdots, \lambda_n
\end{align*}

\paragraph*{Theorem}
\begin{enumerate}
    \item $A$ is diagonalizable $\iff$ it has eigenvectors $\vec{x_1}, \vec{x_2}, \cdots, \vec{x_n}$ and $P = [\vec{x_1} \vec{x_2} \cdots \vec{x_n}]$ is invertible
    \item When this is the case, $P^{-1}AP = $ diag$(\lambda_1, \lambda_2, \cdots, \lambda_n)$ where $\lambda_i$ is the eigenvalue corresponding to $\vec{x_i}$ for $i = 1, 2, \cdots, n$
    \item If $A$ has n \emph{distinct} eigenvalues, then $A$ is diagonalizable. (Solve $(A - \lambda_i I)\vec{x} = \vec{0}$ for each eigenvalue $\lambda_i$ to get n linearly independent eigenvectors.)
\end{enumerate}

\paragraph*{Example}
Let $A = \begin{bmatrix}
    2 & 0 & 0 \\
    1 & 2 & -1 \\
    1 & 3 & -2
\end{bmatrix}$ and $\lambda_1 = 2, lambda_2 = 1, \lambda_3 = -1$. Find a basic eigenvector corresponding to each eigenvalue and determine if $A$ is diagonalizable.
\begin{align*}
    (A - 2I) = &\begin{bmatrix}
    0 & 0 & 0 \\
    1 & 0 & -1 \\
    1 & 3 & -4
    \end{bmatrix} \vec{x} = \vec{0} \implies \vec{x} = t \begin{bmatrix}
    1 \\
    1 \\
    1
    \end{bmatrix} \\
    \therefore \quad &\begin{bmatrix}
    1 \\
    1 \\
    1
    \end{bmatrix} \text{ is a basic eigenvector corresponding to } \lambda_1 = 2 \\
    (A - I) = &\begin{bmatrix}
    1 & 0 & 0 \\
    1 & 1 & -1 \\
    1 & 3 & -1
    \end{bmatrix} \vec{x} = \vec{0} \implies \vec{x} = t \begin{bmatrix}
    0 \\
    1 \\
    1
    \end{bmatrix} \\
    \therefore \quad &\begin{bmatrix}
    0 \\
    1 \\
    1 
    \end{bmatrix} \text{ is a basic eigenvector corresponding to } \lambda_2 = 1 \\
    (A + I) = &\begin{bmatrix}
    3 & 0 & 0 \\
    1 & 3 & -1 \\
    1 & 3 & -1
    \end{bmatrix} \vec{x} = \vec{0} \implies \vec{x} = t \begin{bmatrix}
    0 \\
    1 \\
    3
    \end{bmatrix} \\
    \therefore \quad &\begin{bmatrix}
    0 \\
    1 \\
    3
    \end{bmatrix} \text{ is a basic eigenvector corresponding to } \lambda_3 = -1 \\
    \implies &P = \begin{bmatrix}
    1 & 0 & 0 \\
    1 & 1 & 1 \\
    1 & 1 & 4
    \end{bmatrix} \text{ and} D = \begin{bmatrix}
    2 & 0 & 0 \\
    0 & 1 & 0 \\
    0 & 0 & -1
    \end{bmatrix} \\
\end{align*}

\paragraph*{Example}
Diagonalize $A = \begin{bmatrix}
    0 & 1 & 1 \\
    1 & 0 & 1 \\
    1 & 1 & 0
\end{bmatrix}$
\begin{enumerate}
    \item Charateristic polynomial: $\det(A - \lambda I) = \det\begin{vmatrix}
    -\lambda & 1 & 1 \\
    1 & -\lambda & 1 \\
    1 & 1 & -\lambda
\end{vmatrix} = -\lambda^3 + 3\lambda + 2 = -(\lambda - 2)(\lambda + 1)^2$
    \item Eigenvalues: $\lambda_1 = 2, \lambda_2 = -1$
    \item Eigenvectors: \begin{align*}
        (A - 2I) = &\begin{bmatrix}
        -2 & 1 & 1 \\
        1 & -2 & 1 \\
        1 & 1 & -2
        \end{bmatrix} \vec{x} = \vec{0} \implies \vec{x} = t \begin{bmatrix}
        1 \\
        1 \\
        1
        \end{bmatrix} \\
        (A + I) = &\begin{bmatrix}
        1 & 1 & 1 \\
        1 & 1 & 1 \\
        1 & 1 & 1
        \end{bmatrix} \vec{x} = \vec{0} \implies x_1 + x_2 + x_3 = 0 \implies \vec{x} = t_1 \begin{bmatrix}
        1 \\
        -1 \\
        0
        \end{bmatrix} + t_2 \begin{bmatrix}
        1 \\
        0 \\
        -1
        \end{bmatrix} \\
        \therefore \quad &D = \text{diag}(-1, -1, 2) \text{ and } P = \begin{bmatrix}
            -1 & -1 & 1 \\
            1 & 0 & 1 \\
            0 & 1 & 1
        \end{bmatrix} \text{ is the diagonalized form of } A
    \end{align*}
\end{enumerate}

\subsubsection*{Multiplicity of Eigenvalues}
\paragraph*{Definition}
An eighenvalue $\lambda$ of $A_{n \times n}$ has multiplicity $k$ if it occurs $k$ times as a root of the characteristic polynomial of $A$.
For example: $P_A(\lambda) = (\lambda - 1)^2(\lambda + 2)^3$ has eigenvalues $\lambda_1 = 1$ with multiplicity 2 and $\lambda_2 = -2$ with multiplicity 3.

\paragraph*{Theorem}
An $n \times n$ matrix $A$ is diagonalizable if and only if every eigenvalue $\lambda_i$ of multiplicity $m$ yeilds \emph{exactly} $m$ basic eigenvectors.
In particular, if $A_{n \times n}$ has $n$ distinct eigenvalues, then $A$ is diagonalizable.


\paragraph*{Diagonalization Algorithm}
\begin{enumerate}
    \item Find the eigenvalues of $A$ by solving $\det(A - \lambda I) = 0$
    \item Compute the basic eigenvectors corresponding to each eigenvalue $\lambda_i$ by solving $(A - \lambda_i I)\vec{x} = \vec{0}$
    \item $A$ is diagonalizable $\iff$ there are $n$ basic eigenvectors total. Exactly the multiplicity amount corresponding to each eigenvalue.
    \item In this case, the matrix $P$ is formed with these $n$ \emph{basic} eigenvectors as columns and $D = \text{diag}(\lambda_1, \lambda_2, \cdots, \lambda_n)$ where $\lambda_i$ is the eigenvalue corresponding to the $i$th column of $P$ for $i = 1, 2, \cdots, n$.
\end{enumerate}


\paragraph*{Example}
Show that $A = \begin{bmatrix}
    1 & 1 \\
    0 & 1
\end{bmatrix}$ is not diagonalizable.
\begin{align*}
    &\lambda = 1 \text{ with multiplicity } 2 \implies (A - I) = \begin{bmatrix}
    0 & 1 \\
    0 & 0
\end{bmatrix} \vec{x} = \vec{0} \implies \vec{x} = t \begin{bmatrix}
    1 \\
    0
\end{bmatrix} \\
\therefore \quad &\text{There is only 1 basic eigenvector corresponding to the eigenvalue } 1 \text{ with multiplicity } 2 \\
\implies &A \text{ is not diagonalizable}
\end{align*}

\paragraph*{Example}
Find the diagonalization of $A = \begin{bmatrix}
    1 & 4 & 0 \\
    1 & 1 & 0 \\
    0 & 0 & 3
\end{bmatrix}$.
\begin{enumerate}
    \item Find eigenvalues: 
    \begin{align*}
        \det(A - \lambda I) = &\det\begin{vmatrix}
            1 - \lambda & 4 & 0 \\
            1 & 1 - \lambda & 0 \\
            0 & 0 & 3 - \lambda
        \end{vmatrix} = (3 - \lambda)\det\begin{vmatrix}
            1 - \lambda & 4 \\
            1 & 1 - \lambda
        \end{vmatrix} = (3 - \lambda)((1 - \lambda)^2 - 4) \\
        = &(3 - \lambda)(\lambda^2 - 2\lambda - 3) = (3 - \lambda)(\lambda - 3)(\lambda + 1) = -(3 - \lambda)^2(\lambda + 1)
    \end{align*}
    \item Find basic eigenvectors: 
    \begin{align*}
            (A - 3I) = &\begin{bmatrix}
            -2 & 4 & 0 \\
            1 & -2 & 0 \\
            0 & 0 & 0
        \end{bmatrix} \vec{x} = \vec{0} \implies \vec{x} = t \begin{bmatrix}
            2 \\
            1 \\
            0
        \end{bmatrix} \\
        (A + I) = &\begin{bmatrix}
            0 & 4 & 0 \\
            1 & 0 & 0 \\
            0 & 0 & 4
        \end{bmatrix} \vec{x} = \vec{0} \implies \vec{x} = t \begin{bmatrix}
            0 \\
            0 \\
            1
        \end{bmatrix} \\
        \therefore \quad &D = \text{diag}(3, 3, -1) \text{ and } P = \begin{bmatrix}
            2 & 1 & 0 \\
            1 & 0 & 0 \\
            0 & 0 & 1
        \end{bmatrix} \text{ is the diagonalized form of } A
    \end{align*}
    \item Find $A^3$ using the diagonalization of $A$:
    \begin{align*}
        A^3 = &PD^3P^{-1} = P\begin{bmatrix}
            27 & 0 & 0 \\
            0 & 27 & 0 \\
            0 & 0 & -1
        \end{bmatrix}P^{-1} = \begin{bmatrix}
            2 & 1 & 0 \\
            1 & 0 & 0 \\
            0 & 0 & 1
        \end{bmatrix}\begin{bmatrix}
            27 & 0 & 0 \\
            0 & 27 & 0 \\
            0 & 0 & -1
        \end{bmatrix}\begin{bmatrix}
            0 & 1 & 0 \\
            \frac{1}{2} & -\frac{1}{2} & 0 \\
            0 & 0 & 1
        \end{bmatrix} = \begin{bmatrix}
            54 & -54 & 0 \\
            27 & -27 & 0 \\
            0 & 0 & -1
        \end{bmatrix}
    \end{align*}
\end{enumerate}


\pagebreak

\subsection*{Basic Intro to Jordan Cononical Form}
\paragraph*{Definition}
A matrix $A$ is in Jordan cononical form if it is a block diagonal matrix where each block is a Jordan block. A Jordan block is a square matrix of the form
\begin{equation*}
    J_k(\lambda) = \begin{bmatrix}
        \lambda & 1 & 0 & \cdots & 0 \\
        0 & \lambda & 1 & \cdots & 0 \\
        0 & 0 & \lambda & \cdots & 0 \\
        \vdots & \vdots & \vdots & \ddots & 1 \\
        0 & 0 & 0 & \cdots & \lambda
    \end{bmatrix}_{k \times k}
\end{equation*}

\paragraph*{Theorem}
$J_k(\lambda)$ is not diagonalizable for $k > 1$ because it has only one eigenvalue $\lambda$ with multiplicity $k$ but only one basic eigenvector.

\paragraph*{Theorem}
$\forall A \sim J, J = \text{diag}(J_{k_1}(\lambda_1), J_{k_2}(\lambda_2), \cdots, J_{k_m}(\lambda_m))$ where $J_{k_i}(\lambda_i)$ is a Jordan block corresponding to the eigenvalue $\lambda_i$ for $i = 1, 2, \cdots, m$.

\paragraph*{Example}
$J = \text{diag}(J_1(5), J_2(7), J_1(7), J_3(8))$, find the dimensions of $J$
\begin{equation*}
    J \in \mathbb{R}^{n \times n} \implies n = 1 + 2 + 1 + 3 = 7
\end{equation*}

\paragraph*{Theorem}
$P^{-1}AP = J = \text{diag}(J_{k_1}(\lambda_1), J_{k_2}(\lambda_2), \cdots, J_{k_m}(\lambda_m))$ where $J_{k_i}(\lambda_i)$ is called a block diagonal.\\
$A$ is diagonalizable $\iff$ every Jordan block has size 1. Otherwise, if there is a Jordan block of $P^{-1}AP$ that is larger than $1 \times 1$, then $A$ is not diagonalizable.

\paragraph*{Example}
Show $J_3(\lambda) = \begin{bmatrix}
    \lambda & 1 & 0 \\
    0 & \lambda & 1 \\
    0 & 0 & \lambda
\end{bmatrix}$ is not diagonalizable.
\begin{enumerate}
    \item $J_3(\lambda)$ has only one eigenvalue $\lambda$ with multiplicity 3
    \item $A - \lambda I = \begin{bmatrix}
    0 & 1 & 0 \\
    0 & 0 & 1 \\
    0 & 0 & 0
\end{bmatrix} \vec{x} = \vec{0} \implies \vec{x} = t \begin{bmatrix}
    1 \\
    0 \\
    0
\end{bmatrix}$
    \item There is only one basic eigenvector corresponding to the eigenvalue $\lambda$ with multiplicity 3 $\implies$ $J_3(\lambda)$ is not diagonalizable
\end{enumerate}